% Unit 095 terjemahan Bahasa Indonesia dari chapter7.tex baris 1445--1660.
% Sumber: Methods of Algebra, Volume 2, commit
% 9a5803ff77dd3257484cb177f851a73770a59dd3 (CC BY 4.0).
% stable-unit-id: o014.aljabr2.chapter7.beck-monadicity-b
% Cakupan: sisa Bagian 7.6, dari funktor pembanding hingga versi linear
% Teorema Beck; Unit 096 memulai Bagian 7.7 pada baris 1661.

\hypertarget{o014.aljabr2.chapter7.beck-monadicity-b}{ }
% segment-id: o014.aljabr2.chapter7.beck-monadicity-b.l001
\begin{lemma}\label{prop:T-comparison}
	Tinjau monad $T$ dan komonad $L$ yang berasal dari pasangan adjoin
	$(F,G)$. Misalkan $N\in\Obj(\mathcal{D})$. Data
	% segment-id: o014.aljabr2.chapter7.beck-monadicity-b.q001
	\[ M:=G(N)\in\Obj(\mathcal{C}), \quad
	 a:T(M)=GFG(N)\xrightarrow{G\varepsilon_N}G(N)=M \]
	memberikan $(M,a)\in\Obj(\mathcal{C}^T)$. Dengan demikian diperoleh funktor
	$\mathbb{K}:\mathcal{D}\to\mathcal{C}^T$ yang memenuhi
	$G=U^T\mathbb{K}$.

	Secara dual, $F$ juga terfaktorkan secara kanonik sebagai
	$\mathcal{C}\xrightarrow{\mathbb{K}}\mathcal{D}^L\to\mathcal{D}$, yang
	selanjutnya ditulis $F=U^L\mathbb{K}$, dengan $U^L$ funktor
	pelupa\footnote{Pemakaian lambang $\mathbb{K}$ yang sama tidak akan banyak
	menimbulkan kerancuan, sebab buku ini tidak akan membahas monad dan komonad
	secara bersamaan.}.
\end{lemma}
% segment-id: o014.aljabr2.chapter7.beck-monadicity-b.pf001
\begin{proof}
	Cukup dibahas kasus $T$. Kedua diagram dalam Definisi~\ref{def:T-Alg}
	dalam hal ini menjadi
	% segment-id: o014.aljabr2.chapter7.beck-monadicity-b.g001
	\[\begin{tikzcd}
		GFGFG(N) \arrow[r, "GFG\varepsilon" inner sep=0.5em] \arrow[d, "{G\varepsilon FG}"'] & GFG(N) \arrow[d, "G\varepsilon"] \\
		GFG(N) \arrow[r, "G\varepsilon"'] & G(N)
	\end{tikzcd} \quad \begin{tikzcd}
		G(N) \arrow[r, "\eta G"] \arrow[equal, rd] & GFG(N) \arrow[d, "{G\varepsilon}"] \\
		& G(N).
	\end{tikzcd} \]
	Verifikasi komutativitasnya persis sama dengan yang dilakukan dalam
	Contoh~\ref{eg:adjunction-monad}, sehingga tidak perlu diulangi;
	funktorialitas $N\mapsto(M,a)$ juga langsung terlihat.
\end{proof}

% segment-id: o014.aljabr2.chapter7.beck-monadicity-b.p001
Untuk pasangan adjoin $(F,G)$ yang diberikan, menerapkan funktor merupakan
proses yang melupakan struktur. Lema~\ref{prop:T-comparison} memfaktorkan $G$
(atau $F$) sebagai $U^T\mathbb{K}$ (atau $U^L\mathbb{K}$). Kita ingin
mengetahui apakah struktur hanya hilang pada langkah kedua, yaitu saat
menerapkan funktor pelupa $U^T$. Jika demikian, informasi yang hilang dalam
proses itu dapat direkonstruksi dengan bantuan monad $T$ (atau komonad $L$).
Pertimbangan ini memotivasi konsep berikut.

% segment-id: o014.aljabr2.chapter7.beck-monadicity-b.d001
\begin{definition}\label{def:monadic}
	Tinjau sepasang funktor adjoin
	% segment-id: o014.aljabr2.chapter7.beck-monadicity-b.g002
	$\begin{tikzcd}
		F: \mathcal{C} \arrow[shift left, r] & \mathcal{D}: G \arrow[shift left, l]
	\end{tikzcd}$,
	yang menentukan monad $T$ di atas $\mathcal{C}$ dan komonad $L$ di atas
	$\mathcal{D}$.
	\begin{itemize}
		\item Jika funktor $\mathbb{K}:\mathcal{D}\to\mathcal{C}^T$ dalam
		Lema~\ref{prop:T-comparison} merupakan ekuivalensi kategori, pasangan
		adjoin tersebut disebut \emph{monadik}.
		\item Jika versi dualnya
		$\mathbb{K}:\mathcal{C}\to\mathcal{D}^L$ merupakan ekuivalensi kategori,
		pasangan adjoin tersebut disebut \emph{komonadik}.
	\end{itemize}
\end{definition}

% segment-id: o014.aljabr2.chapter7.beck-monadicity-b.p002
Sebagai contoh, pasangan adjoin bebas--pelupa dalam
Contoh~\ref{eg:Mon-monad} dan \ref{eg:Mod-monad} semuanya monadik; pemeriksaan
langsung bahkan menunjukkan bahwa $\mathbb{K}$ dalam kedua contoh tersebut
merupakan isomorfisme kategori.

% segment-id: o014.aljabr2.chapter7.beck-monadicity-b.p003
Teorema Beck yang akan dibahas berikutnya, juga disebut teorema monadisitas atau
teorema Barr--Beck, mencirikan sifat monadik. Kita memerlukan beberapa
persiapan.

% segment-id: o014.aljabr2.chapter7.beck-monadicity-b.dp001
\begin{definition-proposition}
	\index{fork terbelah@fork terbelah (split fork)}
	Tinjau diagram dalam suatu kategori $\mathcal{C}$
	% segment-id: o014.aljabr2.chapter7.beck-monadicity-b.e001
	\begin{equation*}\label{eqn:split-coeq}\begin{tikzcd}[column sep=large]
		A \arrow[shift left, r, "u"] \arrow[shift right, r, "v" description] & B \arrow[r, "h"] \arrow[dashed, l, bend left, "t"] & Z \arrow[dashed, l, bend left, "s"]
	\end{tikzcd}\end{equation*}
	sedemikian sehingga bagian bergaris penuh komutatif, yaitu $hu=hv$, dan
	$h$ serta $v$ masing-masing mempunyai seksi (yakni invers kanan) $s$ dan
	$t$ seperti ditunjukkan oleh garis putus-putus, dengan
	$sh=ut\in\End(B)$. Bagian bergaris penuh dari diagram yang mempunyai sifat
	ini disebut \emph{garpu terbelah}. Diagram tersebut secara otomatis
	memberikan kopenyama dari $u$ dan $v$ dalam $\mathcal{C}$.
\end{definition-proposition}
% segment-id: o014.aljabr2.chapter7.beck-monadicity-b.pf002
\begin{proof}
	Tinjau diagram berikut:
	% segment-id: o014.aljabr2.chapter7.beck-monadicity-b.g003
	\[\begin{tikzcd}[row sep=small]
		& & W \\
		A \arrow[shift left, r, "u"] \arrow[shift right, r, "v"'] & B \arrow[r, "h"'] \arrow[ru, "k"] & Z
	\end{tikzcd} \qquad \begin{array}{ll}
		W \in \Obj(\mathcal{C}) \\
		ku = kv.
	\end{array} \]
	Jika terdapat $\varphi:Z\to W$ dengan $\varphi h=k$, maka
	$\varphi=\varphi hs=ks$. Sebaliknya, $\varphi:=ks$ memang memenuhi
	$\varphi h=ksh=kut=kvt=k$. Ini membuktikan sifat universal yang
	diperlukan.
\end{proof}

% segment-id: o014.aljabr2.chapter7.beck-monadicity-b.p004
Berbeda dengan kopenyama pada umumnya, garpu terbelah bersifat
``mutlak'': citranya di bawah funktor apa pun tetap merupakan garpu terbelah.

% segment-id: o014.aljabr2.chapter7.beck-monadicity-b.eg001
\begin{example}\label{eg:T-split-fork}
	Misalkan $(T,\mu,\eta)$ monad di atas $\mathcal{C}$. Setiap
	$(M,a)\in\Obj(\mathcal{C}^T)$ memberikan garpu terbelah dalam
	$\mathcal{C}$:
	% segment-id: o014.aljabr2.chapter7.beck-monadicity-b.e002
	\begin{equation*}\begin{tikzcd}[column sep=large]
			T^2(M) \arrow[shift left, r, "Ta"] \arrow[shift right, r, "\mu_M" description] & T(M) \arrow[r, "a"] \arrow[dashed, bend left, l, "\eta_{TM}"] & M \arrow[dashed, bend left, l, "\eta_M"]
		\end{tikzcd} \quad \begin{array}{ll}
			a (Ta) = a \mu_M , & \eta_M a = (Ta) \eta_{TM}, \\
			a \eta_M = \identity_M , & \mu_M \eta_{TM} = \identity_{TM}.
		\end{array}\end{equation*}
	Kesamaan-kesamaan di sebelah kanan adalah bagian dari
	Definisi~\ref{def:monad} dan \ref{def:T-Alg}; sebagai contoh,
	$\eta_Ma=(Ta)\eta_{TM}$ berasal dari kealamian
	$\eta:\identity_{\mathcal{C}}\to T$. Secara khusus, diagram di atas
	merealisasikan $M$ sebagai kopenyama dari $Ta$ dan $\mu_M$.
\end{example}

% segment-id: o014.aljabr2.chapter7.beck-monadicity-b.d002
\begin{definition}\label{def:Beck-split-pair}
	Tinjau funktor $G:\mathcal{D}\to\mathcal{C}$ dan sepasang morfisme
	$f,g:X\rightrightarrows Y$ dalam $\mathcal{D}$. Jika terdapat morfisme
	$h:G(Y)\to Z$ sedemikian sehingga diagram
	% segment-id: o014.aljabr2.chapter7.beck-monadicity-b.e003
	\begin{equation*}\begin{tikzcd}
			G(X) \arrow[shift left, r, "Gf"] \arrow[shift right, r, "Gg"'] & G(Y) \arrow[r, "h"] & Z
	\end{tikzcd}\end{equation*}
	merupakan garpu terbelah dalam $\mathcal{C}$, maka $(f,g)$ disebut
	\emph{pasangan $G$-terbelah} dalam $\mathcal{D}$.
\end{definition}

% segment-id: o014.aljabr2.chapter7.beck-monadicity-b.p005
Misalkan $(f,g)$ pasangan $G$-terbelah. Karena $(Gf,Gg)$ mempunyai
kopenyama, wajar ditanyakan apakah kopenyama tersebut dapat diangkat menjadi
kopenyama dari $(f,g)$ dalam $\mathcal{D}$, dan apakah pengangkatan itu unik.
Kopenyama merupakan kasus khusus dari $\varinjlim$; konsep funktor yang
menciptakan atau mempertahankan $\varinjlim$ dalam
Definisi~\ref{def:create-limit} menyediakan istilah yang tepat untuk masalah
ini.

% segment-id: o014.aljabr2.chapter7.beck-monadicity-b.p006
Perhatikan bahwa funktor pelupa $U^T:\mathcal{C}^T\to\mathcal{C}$ merupakan
funktor konservatif dalam pengertian Definisi~\ref{def:conservative}.

% segment-id: o014.aljabr2.chapter7.beck-monadicity-b.l002
\begin{lemma}\label{prop:UT-split-fork}
	Misalkan $(T,\mu,\eta)$ monad di atas $\mathcal{C}$. Funktor $U^T$
	menciptakan kopenyama dari setiap pasangan $U^T$-terbelah.
\end{lemma}
% segment-id: o014.aljabr2.chapter7.beck-monadicity-b.pf003
\begin{proof}
	Misalkan diberikan sepasang morfisme dalam $\mathcal{C}^T$,
	$u,v:(M,a)\rightrightarrows(M',a')$, dan sebuah garpu terbelah dalam
	$\mathcal{C}$:
	% segment-id: o014.aljabr2.chapter7.beck-monadicity-b.e004
	\begin{equation*}\begin{tikzcd}[column sep=large]
		M \arrow[shift left, r, "u"] \arrow[shift right, r, "v" description] & M' \arrow[r, "h"] \arrow[dashed, l, bend left, "t"] & M'' \arrow[dashed, l, bend left, "s"].
	\end{tikzcd}\end{equation*}
	Karena diagram ini memberikan kopenyama dalam $\mathcal{C}$, cukup
	memperluas $M''$ secara unik menjadi
	$(M'',a'')\in\Obj(\mathcal{C}^T)$ sedemikian sehingga $h$ merupakan
	morfisme dalam $\mathcal{C}^T$, kemudian menunjukkan bahwa konstruksi ini
	memberikan kopenyama dari $u$ dan $v$ dalam $\mathcal{C}^T$.

	Citra garpu terbelah di bawah funktor apa pun tetap merupakan garpu
	terbelah. Karena itu, dalam diagram komutatif pada bagian bergaris penuh
	% segment-id: o014.aljabr2.chapter7.beck-monadicity-b.g004
	\[\begin{tikzcd}
		T(M) \arrow[shift left, r, "Tu"] \arrow[shift right, r, "Tv"'] \arrow[d, "a"'] & T(M') \arrow[r, "Th"] \arrow[d, "{a'}"] & T(M'') \arrow[dashed, d, "{a''}"] \\
		M \arrow[shift left, r, "u"] \arrow[shift right, r, "v"'] & M' \arrow[r, "h"'] & M''
	\end{tikzcd}\]
	kedua baris merupakan kopenyama dalam $\mathcal{C}$. Sifat universal
	karenanya memberikan tepat satu $a''$ yang membuat seluruh diagram
	komutatif. Jika $(M'',a'')$ terbukti merupakan aljabar-$T$, maka
	komutativitas bagian kanan diagram itu tepat berarti bahwa $h$ merupakan
	morfisme dalam $\mathcal{C}^T$.

	Untuk itu, yang perlu diverifikasi adalah komutativitas kedua diagram
	berikut.\footnote{Catatan penerjemah: simpul kanan bawah diagram pertama
	tercetak sebagai $T(M'')$ dalam sumber. Hukum satuan aljabar-$T$ mengharuskan
	simpul itu berupa $M''$; target tersebut dikoreksi di sini.}
	% segment-id: o014.aljabr2.chapter7.beck-monadicity-b.g005
	\[\begin{tikzcd}
		M'' \arrow[r, "{\eta_{M''}}"] \arrow[equal, rd] & T(M'') \arrow[d, "{a''}"] \\
		& M''
	\end{tikzcd} \quad
	\begin{tikzcd}
		T^2(M'') \arrow[r, "{\mu_{M''}}"] \arrow[d, "{Ta''}"'] & T(M'') \arrow[d, "{a''}"] \\
		T(M'') \arrow[r, "{a''}"'] & M'' .
	\end{tikzcd}\]
	Diagram kopenyama menunjukkan bahwa $h$ dan $Th$ adalah epimorfisme
	(demikian pula $T^2h$). Oleh karena itu, komutativitas kedua diagram di atas
	langsung mengikuti dari komutativitas diagram yang bersesuaian untuk $M$
	dan $M'$.

	Akhirnya, akan ditunjukkan bahwa $h$ memberikan kopenyama dari $u$ dan $v$
	dalam $\mathcal{C}^T$. Misalkan terdapat morfisme dalam $\mathcal{C}^T$
	$k:(M',a')\to(W,b)$ dengan $ku=kv$.\footnote{Catatan penerjemah: sumber
	mencetak domain $k$ sebagai $(M'',a'')$, yang membuat $ku$, $kv$, dan
	$\varphi h=k$ tidak terdefinisi. Domain yang benar adalah $(M',a')$.}
	Dalam $\mathcal{C}$ terdapat morfisme tunggal $\varphi:M''\to W$ dengan
	$\varphi h=k$. Tinggal dibuktikan bahwa $\varphi$ sebenarnya merupakan
	morfisme dalam $\mathcal{C}^T$. Tinjau diagram
	% segment-id: o014.aljabr2.chapter7.beck-monadicity-b.g006
	\[\begin{tikzcd}
		T(M') \arrow[r, "Th"] \arrow[d, "{a'}"'] & T(M'') \arrow[r, "T\varphi"] \arrow[d, "{a''}"] & T(W) \arrow[d, "b"] \\
		M' \arrow[r, "h"'] & M'' \arrow[r, "\varphi"'] & W.
	\end{tikzcd}\]
	Komposisi masing-masing baris adalah $Tk$ dan $k$. Karena $h$ dan $k$
	merupakan morfisme dalam $\mathcal{C}^T$, seluruh batas luar dan persegi
	kiri komutatif. Maka persegi kanan komutatif setelah dikomposisikan dengan
	$Th$. Karena $Th$ epimorfisme, persegi kanan itu sendiri komutatif. Hal ini
	menunjukkan bahwa $\varphi$ merupakan morfisme dalam $\mathcal{C}^T$.
\end{proof}

% segment-id: o014.aljabr2.chapter7.beck-monadicity-b.l003
\begin{lemma}\label{prop:Beck-FG-coeq}
	Misalkan funktor $G:\mathcal{D}\to\mathcal{C}$ mempunyai adjoin kiri $F$
	dan $G$ menciptakan kopenyama yang bersesuaian untuk semua pasangan
	$G$-terbelah $f,g:X\rightrightarrows Y$
	(Definisi~\ref{def:Beck-split-pair}). Maka diagram berikut merupakan
	kopenyama:
	% segment-id: o014.aljabr2.chapter7.beck-monadicity-b.e005
	\begin{equation*}\begin{tikzcd}[column sep=large]
		FGFG(N) \arrow[shift left, r, "FG\varepsilon_N"] \arrow[shift right, r, "\varepsilon_{FG(N)}"'] & FG(N) \arrow[r, "\varepsilon_N"] & N,
	\end{tikzcd}\end{equation*}
	dengan $N\in\Obj(\mathcal{D})$.
\end{lemma}
% segment-id: o014.aljabr2.chapter7.beck-monadicity-b.pf004
\begin{proof}
	Menerapkan $G$ pada diagram tersebut menghasilkan garpu terbelah dari
	Contoh~\ref{eg:T-split-fork} untuk
	$(M,a):=(G(N),G\varepsilon_N)=\mathbb{K}(N)$. Karena $G$ menciptakan
	kopenyama yang bersesuaian, diagram semula merupakan kopenyama.
\end{proof}

% segment-id: o014.aljabr2.chapter7.beck-monadicity-b.t001
\begin{theorem}[J.\ M.\ Beck]\label{prop:Beck}
	\index{Teorema Beck@Teorema Beck}
	Misalkan funktor $G:\mathcal{D}\to\mathcal{C}$ mempunyai adjoin kiri $F$.
	Pernyataan-pernyataan berikut ekuivalen:
	\begin{enumerate}[(i)]
		\item pasangan adjoin $(F,G)$ bersifat monadik
		(Definisi~\ref{def:monadic});
		\item $G$ menciptakan kopenyama yang bersesuaian untuk setiap pasangan
		$G$-terbelah $f,g:X\rightrightarrows Y$;
		\item $G$ konservatif (Definisi~\ref{def:conservative}), setiap pasangan
		$G$-terbelah $f,g:X\rightrightarrows Y$ mempunyai kopenyama
		$X\rightrightarrows Y\to Z$ dalam $\mathcal{D}$, dan citra kopenyama
		tersebut di bawah $G$ merupakan kopenyama dari $(Gf,Gg)$.
	\end{enumerate}

	Jika salah satu kondisi di atas berlaku, suatu funktor kuasi-invers
	$\mathbb{L}$ bagi $\mathbb{K}:\mathcal{D}\to\mathcal{C}^T$ dapat
	dideskripsikan sebagai berikut. Untuk $(M,a)\in\Obj(\mathcal{C}^T)$ terdapat
	diagram kopenyama
	% segment-id: o014.aljabr2.chapter7.beck-monadicity-b.e006
	\begin{equation}\label{eqn:Beck-aux1}\begin{tikzcd}
		FGF(M) \arrow[shift left, r, "Fa"] \arrow[shift right, r, "\varepsilon_{FM}"'] & F(M) \arrow[r] & \mathbb{L}(M, a).
	\end{tikzcd}\end{equation}

	Untuk komonad yang ditentukan oleh $(F,G)$, kriterianya sepenuhnya dual.
\end{theorem}
% segment-id: o014.aljabr2.chapter7.beck-monadicity-b.pf005
\begin{proof}
	Pertama, kita buktikan (i) $\implies$ (ii). Misalkan
	$\mathbb{K}:\mathcal{D}\to\mathcal{C}^T$ mempunyai funktor kuasi-invers
	$\mathbb{L}$. Sifat-sifat berupa pasangan terbelah, garpu terbelah, dan
	penciptaan $\varinjlim$ tidak berubah di bawah ekuivalensi kategori. Karena
	$G=U^T\mathbb{K}$, cukup dibuktikan bahwa
	$U^T:\mathcal{C}^T\to\mathcal{C}$ menciptakan kopenyama yang bersesuaian
	untuk setiap pasangan $U^T$-terbelah. Ini tepat isi
	Lema~\ref{prop:UT-split-fork}.

	Selanjutnya, (i) $\implies$ (iii). Karena (i) $\implies$ (ii) telah
	diketahui, tinggal ditunjukkan bahwa $G$ konservatif apabila $G$ monadik.
	Akan tetapi, $G=U^T\mathbb{K}$, sedangkan $\mathbb{K}$ merupakan
	ekuivalensi kategori dan $U^T$ jelas konservatif. Jadi, $G$ juga
	konservatif.

	Sekarang kita buktikan (ii) $\implies$ (i) dengan membangun funktor kuasi-invers
	$\mathbb{L}$ bagi $\mathbb{K}$. Untuk
	$(M,a)\in\Obj(\mathcal{C}^T)$, tinjau pasangan morfisme dalam
	$\mathcal{D}$
	% segment-id: o014.aljabr2.chapter7.beck-monadicity-b.g007
	$\begin{tikzcd}
		FGF(M) \arrow[shift left, r, "Fa"] \arrow[shift right, r, "\varepsilon_{FM}"'] & F(M).
	\end{tikzcd}$
	Citranya di bawah $G$ adalah
	% segment-id: o014.aljabr2.chapter7.beck-monadicity-b.g008
	$\begin{tikzcd}
		T^2(M) \arrow[shift left, r, "Ta"] \arrow[shift right, r, "\mu_M"'] & T(M)
	\end{tikzcd}$\!,
	yang menurut Contoh~\ref{eg:T-split-fork} dapat diperluas menjadi garpu
	terbelah. Jadi, $(Fa,\varepsilon_{FM})$ merupakan pasangan $G$-terbelah,
	dan diagram semula dapat diperluas menjadi diagram kopenyama dalam
	$\mathcal{D}$, yaitu tepat \eqref{eqn:Beck-aux1}.

	Dari \eqref{eqn:Beck-aux1} dan sifat universal kopenyama, jika diberikan
	morfisme $(M,a)\xrightarrow{\varphi}(M',a')$ dalam $\mathcal{C}^T$, maka
	terdapat morfisme tunggal
	$\mathbb{L}(M,a)\to\mathbb{L}(M',a')$ yang membuat diagram berikut
	komutatif:
	% segment-id: o014.aljabr2.chapter7.beck-monadicity-b.e007
	\begin{equation*}\begin{tikzcd}
		FGF(M) \arrow[shift left, r, "Fa"] \arrow[shift right, r, "\varepsilon_{FM}"'] \arrow[d, "FGF\varphi"'] & F(M) \arrow[r] \arrow[d, "F\varphi"] & \mathbb{L}(M, a) \arrow[d] \\
		FGF(M') \arrow[shift left, r, "{Fa'}"] \arrow[shift right, r, "\varepsilon_{FM'}"'] & F(M') \arrow[r] & \mathbb{L}(M', a').
	\end{tikzcd}\end{equation*}
	Dengan demikian, $(M,a)\mapsto\mathbb{L}(M,a)$ menjadi funktor
	$\mathbb{L}:\mathcal{C}^T\to\mathcal{D}$.

	Berikutnya akan dibuktikan
	$\mathbb{L}\mathbb{K}\simeq\identity_{\mathcal{D}}$.\footnote{Catatan
	penerjemah: sumber mencetak $\identity_{\mathcal{C}}$ di sini dan pada
	kalimat pembanding berikutnya. Karena $\mathbb{L}\mathbb{K}$ adalah
	endofunktor $\mathcal{D}$, subskrip yang benar adalah $\mathcal{D}$.}
	Untuk setiap $N\in\Obj(\mathcal{D})$, konstruksi di atas memberikan diagram
	kopenyama
	% segment-id: o014.aljabr2.chapter7.beck-monadicity-b.e008
	\begin{equation*}\begin{tikzcd}[column sep=large]
		FGFG(N) \arrow[shift left, r, "FG\varepsilon_N"] \arrow[shift right, r, "\varepsilon_{FG(N)}"'] & FG(N) \arrow[r] & \mathbb{L}\mathbb{K}(N).
	\end{tikzcd}\end{equation*}

	Membandingkannya dengan diagram kopenyama dalam
	Lema~\ref{prop:Beck-FG-coeq} langsung memberikan isomorfisme kanonik
	$\mathbb{L}\mathbb{K}\simeq\identity_{\mathcal{D}}$.

	Selanjutnya, kita verifikasi
	$\mathbb{K}\mathbb{L}\simeq\identity_{\mathcal{C}^T}$. Jelas bahwa
	$\mathbb{K}F=\mathrm{Free}^T$; keduanya memetakan objek $M$ ke
	$(GF(M),G\varepsilon_{FM})$. Menurut definisi dan
	$\mu:=G\varepsilon F$, menerapkan $\mathbb{K}$ pada
	\eqref{eqn:Beck-aux1} menghasilkan
	% segment-id: o014.aljabr2.chapter7.beck-monadicity-b.e009
	\begin{equation}\label{eqn:Beck-aux3}
		\begin{tikzcd}[row sep=small]
			(T^2(M), \ldots) \arrow[shift left, r, "Ta"] \arrow[shift right, r, "\mu_M"'] & (T(M), \ldots) \arrow[r] & \mathbb{K}\mathbb{L}(M, a). \\
			\mathrm{Free}^T(T(M)) \arrow[equal, u] & \mathrm{Free}^T(M) \arrow[equal, u] &
		\end{tikzcd}
	\end{equation}

	Menerapkan funktor pelupa $U^T$ pada dua suku di sebelah kiri
	\eqref{eqn:Beck-aux3} menghasilkan
	% segment-id: o014.aljabr2.chapter7.beck-monadicity-b.g009
	$\begin{tikzcd}
		T^2(M) \arrow[shift left, r, "{Ta}"] \arrow[shift right, r, "{\mu_M}"'] & T(M)
	\end{tikzcd}$.
	Contoh~\ref{eg:T-split-fork} memperluas diagram ini menjadi garpu terbelah
	dengan kopenyama $M$. Di sisi lain, $U^T$ menciptakan kopenyama ini
	(Lema~\ref{prop:UT-split-fork}). Dengan mengingat
	Definisi~\ref{def:create-limit}, diperoleh bahwa \eqref{eqn:Beck-aux3} juga
	merupakan diagram kopenyama.

	Di pihak lain, terapkan Lema~\ref{prop:Beck-FG-coeq} pada pasangan adjoin
	$(\mathrm{Free}^T,U^T)$ dan objek $N=(M,a)$. Dengan menguraikan deskripsi
	pasangan adjoin ini secara saksama, diperoleh diagram kopenyama
	% segment-id: o014.aljabr2.chapter7.beck-monadicity-b.g010
	\[\begin{tikzcd}[row sep=tiny]
		\mathrm{Free}^T(T(M)) \arrow[shift left, r, "Ta"] \arrow[shift right, r, "\mu_M"'] & \mathrm{Free}^T(M) \arrow[r] & (M, a).
	\end{tikzcd}\]
	Perbandingan dengan paragraf sebelumnya langsung memberikan
	$\mathbb{K}\mathbb{L}\simeq\identity_{\mathcal{C}^T}$.

	Terakhir, kita buktikan (iii) $\implies$ (ii). Bagian kedua dari (iii) tepat
	menyatakan bahwa semua pasangan $G$-terbelah mempunyai kopenyama dan $G$
	mempertahankan kopenyama tersebut. Karena $G$ konservatif,
	Catatan~\ref{rem:conservative-reflect} menunjukkan bahwa (ii) berlaku.

	Dengan membalik semua morfisme tetapi mempertahankan arah funktor, yaitu
	dengan mengganti $\mathcal{C}$ dan $\mathcal{D}$ oleh
	$\mathcal{C}^{\opp}$ dan $\mathcal{D}^{\opp}$, diperoleh karakterisasi
	komonadisitas. Perinciannya dihilangkan.
\end{proof}

% segment-id: o014.aljabr2.chapter7.beck-monadicity-b.co001
\begin{corollary}
	Tinjau sepasang funktor adjoin
	% segment-id: o014.aljabr2.chapter7.beck-monadicity-b.g011
	$\begin{tikzcd}
		F: \mathcal{C} \arrow[shift left, r] & \mathcal{D}: G \arrow[shift left, l]
	\end{tikzcd}$
	beserta morfisme unit $\eta$ dan kounit $\varepsilon$ yang bersesuaian.
	Jika pasangan adjoin ini monadik, maka diagram
	% segment-id: o014.aljabr2.chapter7.beck-monadicity-b.g012
	\[\begin{tikzcd}[column sep=large]
		FGFG(N) \arrow[shift left, r, "FG\varepsilon_N"] \arrow[shift right, r, "\varepsilon_{FG(N)}"'] & FG(N) \arrow[r, "\varepsilon_N"] & N
	\end{tikzcd}\]
	merupakan kopenyama untuk setiap $N\in\Obj(\mathcal{D})$.
\end{corollary}
% segment-id: o014.aljabr2.chapter7.beck-monadicity-b.pf006
\begin{proof}
	Gabungkan Lema~\ref{prop:Beck-FG-coeq} dengan
	Teorema~\ref{prop:Beck}~(ii).
\end{proof}

% segment-id: o014.aljabr2.chapter7.beck-monadicity-b.p007
Dalam praktik, versi linear Teorema Beck sering diperlukan. Lebih tepatnya,
misalkan $\Bbbk$ gelanggang komutatif dan $(F,G)$ pasangan adjoin antara
kategori-$\Bbbk\dcate{Mod}$, dan kedua funktor tersebut $\Bbbk$-linear;
lihat Definisi~\ref{def:cat-k-linear} dan
Proposisi~\ref{prop:automatic-k-morphism}. Mudah dilihat bahwa
$\mathcal{C}^T$ secara alami menjadi kategori-$\Bbbk\dcate{Mod}$ dan, menurut
definisi, kedua funktor dalam
$\mathcal{D}\xrightarrow{\mathbb{K}}\mathcal{C}^T\xrightarrow{U^T}\mathcal{C}$
juga $\Bbbk$-linear. Jika kondisi (ii) atau (iii) dalam
Teorema Beck~\ref{prop:Beck} terpenuhi, deskripsi melalui diagram
\eqref{eqn:Beck-aux1} langsung menunjukkan bahwa funktor kuasi-invers
$\mathbb{L}$ bagi $\mathbb{K}$ juga $\Bbbk$-linear. Dengan demikian diperoleh
ekuivalensi $\Bbbk$-linear
$\mathbb{K}:\mathcal{D}\to\mathcal{C}^T$.
